\section{Introdução}
Representar uma curva por segmentos afins envolve erro, número de peças e tamanho aritmético dos parâmetros. Este artigo articula esses recursos em uma biblioteca de segmentos racionais ancorados em células diádicas. Cada átomo possui endereço único e uma interface global que recupera seu suporte.

A aproximação linear por partes, a escolha de nós e a aproximação não linear têm tratamento clássico em \citet{deboor1973,burchard1974,deboor2001,powell1981,schumaker2007,devore1998}. A representação por hinges e sua relação com redes ReLU são desenvolvidas em \citet{yarotsky2017,devore2021,daubechies2022}; a aproximação com pesos codificáveis, em \citet{petersen2018,elbrachter2021,bolcskei2019,guhring2021}. Aqui, esses resultados se integram a uma parametrização racional com recuperação exata da célula e controle conjunto de erro, altura e bits.

A construção usa a altura elementar dos racionais \citep{bombieri2006} e distingue a enumeração por altura da ordem de Calkin--Wilf \citep{calkin2000}, relacionada à sequência de Stern \citep{stern1858,northshield2010}. A análise de descrição emprega códigos universais de inteiros \citep{elias1975,cover2006} e entropia métrica \citep{kolmogorov1959,kolmogorov1961,lorentz1996}. O refinamento combina aproximação em árvores \citep{cohen2001,binev2004,kessebohmer2025} com aritmética intervalar \citep{moore2009,silva2025}. A existência de minimizadores será obtida em catálogos finitos, dentro do contexto de otimização em classes de aproximação \citep{christof2021}.

O modelo fixa inclinações não nulas para tornar o suporte recuperável pela imagem do segmento. As contagens de altura e os limites minimax referem-se ao calibre definido na seção~7. Os comprimentos de serialização pressupõem domínio e convenção de fronteira compartilhados; a certificação adaptativa requer majorantes validados da segunda derivada. Os custos de seleção, armazenamento do decodificador e tempo de execução são recursos definidos separadamente. A quantização diádica também será distinguida da escolha de nós livres.

A exposição progride da indexação à biblioteca, das montagens às garantias de aproximação e destas à codificação. As contribuições são a enumeração invertível de jatos racionais, a contagem dos átomos e o endereçamento minimax sob altura, os formatos prefix-free, o controle simultâneo de erro, altura e descrição e a construção de catálogos finitos com representação canônica.
